\documentclass[final]{article}

% if you need to pass options to natbib, use, e.g.:
%     \PassOptionsToPackage{numbers, compress}{natbib}
% before loading neurips_2025

% The authors should use one of these tracks.
% Before accepting by the NeurIPS conference, select one of the options below.
% 0. "default" for submission
 \usepackage{neurips_2025}
% the "default" option is equal to the "main" option, which is used for the Main Track with double-blind reviewing.
% 1. "main" option is used for the Main Track
%  \usepackage[main]{neurips_2025}
% 2. "position" option is used for the Position Paper Track
%  \usepackage[position]{neurips_2025}
% 3. "dandb" option is used for the Datasets & Benchmarks Track
 % \usepackage[dandb]{neurips_2025}
% 4. "creativeai" option is used for the Creative AI Track
%  \usepackage[creativeai]{neurips_2025}
% 5. "sglblindworkshop" option is used for the Workshop with single-blind reviewing
 % \usepackage[sglblindworkshop]{neurips_2025}
% 6. "dblblindworkshop" option is used for the Workshop with double-blind reviewing
%  \usepackage[dblblindworkshop]{neurips_2025}

% After being accepted, the authors should add "final" behind the track to compile a camera-ready version.
% 1. Main Track
 % \usepackage[main, final]{neurips_2025}
% 2. Position Paper Track
%  \usepackage[position, final]{neurips_2025}
% 3. Datasets & Benchmarks Track
 % \usepackage[dandb, final]{neurips_2025}
% 4. Creative AI Track
%  \usepackage[creativeai, final]{neurips_2025}
% 5. Workshop with single-blind reviewing
%  \usepackage[sglblindworkshop, final]{neurips_2025}
% 6. Workshop with double-blind reviewing
%  \usepackage[dblblindworkshop, final]{neurips_2025}
% Note. For the workshop paper template, both \title{} and \workshoptitle{} are required, with the former indicating the paper title shown in the title and the latter indicating the workshop title displayed in the footnote.
% For workshops (5., 6.), the authors should add the name of the workshop, "\workshoptitle" command is used to set the workshop title.
% \workshoptitle{WORKSHOP TITLE}

% "preprint" option is used for arXiv or other preprint submissions
 % \usepackage[preprint]{neurips_2025}

% to avoid loading the natbib package, add option nonatbib:
%    \usepackage[nonatbib]{neurips_2025}

\usepackage[utf8]{inputenc} % allow utf-8 input
\usepackage[T1]{fontenc}    % use 8-bit T1 fonts
\usepackage{hyperref}       % hyperlinks
\usepackage{url}            % simple URL typesetting
\usepackage{booktabs}       % professional-quality tables
\usepackage{amsfonts}       % blackboard math symbols
\usepackage{nicefrac}       % compact symbols for 1/2, etc.
\usepackage{microtype}      % microtypography
\usepackage{xcolor}         % colors

% Note. For the workshop paper template, both \title{} and \workshoptitle{} are required, with the former indicating the paper title shown in the title and the latter indicating the workshop title displayed in the footnote. 
\title{Rapport \\
        Projet avancé en apprentissage automatique \\
        IFT - 3710}


% The \author macro works with any number of authors. There are two commands
% used to separate the names and addresses of multiple authors: \And and \AND.
%
% Using \And between authors leaves it to LaTeX to determine where to break the
% lines. Using \AND forces a line break at that point. So, if LaTeX puts 3 of 4
% authors names on the first line, and the last on the second line, try using
% \AND instead of \And before the third author name.


\author{
    Cédric Guévremont \\
    Matricule : 20266532 \\
    \texttt{cedric.guevremont@umontreal.ca}
    \And
    Rima Boujenane \\
     Matricule : 20235550 \\
    \texttt{rima.boujenane@umontreal.ca}
    \And
    Abdelghafour Rahmouni \\
    Matricule : 20246224 \\
    \texttt{abdel.ghafour.rahmouni@umontreal.ca}
}


\begin{document}


\maketitle


\section{Description du projet}
    \subsection{Contexte}
        Dans ce projet, on considère un agent d’apprentissage automatique chargé de gérer un portefeuille composé d’un ensemble
        limité d’actifs financiers et un capital de départ. À chaque pas de temps, l’agent observe l’état du marché à partir de données historiques
        (rendements, volatilité, indicateurs techniques) et peut effectuer des actions élémentaires sur ces actifs : acheter,
        vendre ou conserver. Ces actions modifient l’allocation du portefeuille et génèrent une récompense dépendant de
        la performance et du risque.\\
        Ce cadre s’inscrit dans une approche d’apprentissage par renforcement orientée vers une gestion responsable
        du risque, visant à limiter les comportements spéculatifs excessifs, à favoriser des stratégies d’investissement stables
        à long terme, dans une perspective d'usage responsable des méthodes d’apprentissage automatique.

    \subsection{Objectifs}
        L’objectif de ce projet est d’expérimenter l’usage de réseaux récurrents (RNN/LSTM) combiné à l’apprentissage par
        renforcement, afin d’évaluer dans quelle mesure ces approches permettent de réduire le risque global d’un portefeuille
        financier dans un contexte de prise de décision séquentielle.\\\\
        Plus largement, ce travail ouvre la possibilité d’examiner la transférabilité du cadre proposé à d’autres domaines
        reposant sur des séries temporelles et des décisions séquentielles, tels que l’énergie, la logistique ou la gestion
        de ressources.\\\\
        Dans cette perspective, le projet considère notamment comme axes d’exploration la minimisation de la volatilité du
        portefeuille sur l’horizon d’investissement, la réduction des pertes maximales lors des phases
        de stress du marché, l’apprentissage de décisions d’achat, de vente et de conservation cohérentes dans le temps,
        limitant les réallocations excessives, ainsi que l’analyse de l’influence de fonctions de récompense orientées vers
        le contrôle du risque et la stabilité sur les politiques d’allocation apprises.

    \subsection{Méthodes prévues}
        \begin{itemize}
            \item \textbf{Observations} : Fenêtres temporelles de données financières historiques (rendements, volatilité,
                indicateurs techniques) extraites via l’API de Yahoo Finance.

            \item \textbf{Modélisation} : Encodeur d’état basé sur des réseaux récurrents de type LSTM, couplé à un agent
                d’apprentissage par renforcement profond de type Actor-Critic (PPO ou DDPG).

            \item \textbf{Fonction de récompense} : Combinaison du rendement du portefeuille et de pénalités liées au risque
                (volatilité) et aux réallocations excessives.

            \item \textbf{Évaluation} : Backtesting temporel avec comparaison à un indice de référence tel que le S\&P 500.

        \end{itemize}
    \subsection{Résultats prévus}
        Les résultats attendus incluent l’obtention d’un rendement global positif sur l’horizon de test, indiquant que
        l’agent apprend une stratégie de gestion effective plutôt qu’un comportement trivial. On s’attend également à
        ce que la politique apprise produise un rendement relativement stable, caractérisé par une volatilité maîtrisée
        et des pertes limitées lors des périodes de stress. Les performances devraient se situer à un niveau comparable
        à celui du marché, tel que le S\&P 500, tout en offrant un meilleur contrôle du risque, et montrer que l’agent a
        effectivement appris une stratégie d’allocation cohérente plutôt que de simplement rester inactif.

\section{Related works}
    \subsection{Foundational Financial Theories}
    Traditional portfolio optimization methods, such as Markowitz's Modern Portfolio Theory (MPT) \textbf{markowitz1952}
    have been use for decades to balance risk and return.
    However, with the rise of machine learning, new approaches have emerged to tackle portfolio management.
    \textbf{Jang & Seong, 2023} suggested the use of Deep Reinforcement Learning (DRL) combined with Technical Analysis
    and Modern Porfolio Theory to optimize portfolio management.

\end{document}
